% Tự sinh từ rule.csv bởi generate_tex.py; không chỉnh sửa thủ công.
Sau quá trình phân tích, khóa luận đã bổ sung \CustomRulesWritten Custom Rules bao gồm:
\begin{enumerate}
\item \textbf{Local Scheduled Task Creation V2}. Phát hiện hành vi tạo Scheduled Task thông qua event tạo file (Event ID 11) trong thư mục C:\textbackslash{}Windows\textbackslash{}System32\textbackslash{}Tasks\textbackslash{} và C:\textbackslash{}Windows\textbackslash{}Tasks\textbackslash{}*. Ngoài ra, có thể dùng event id 4698, tuy nhiên trường winlog.event\_data.TaskName chứa tên Scheduled Task không nằm trong ECS nên sẽ khó tuning hơn. Rule phát hiện được các biến thể tạo scheduled task khác nhau thay vì chỉ phát hiện schtasks.exe như rule Local Scheduled Task Creation mặc định
\item \textbf{Modify user account}. Phát hiện hành vi enable user thông qua event id 4722 (Enable user) và 4738 (A user account was changed); hành vi tạo mới user thông qua tiến trình net user.
\item \textbf{Suspicious Internet Settings Registry Modification}. Phát hiện các hành vi can thiệp bất thường vào cấu hình Internet Settings trong Registry để hạ thấp cơ chế bảo mật duyệt web hoặc phục vụ C2; bao gồm: thêm tên miền vào Trusted Sites (ZoneMap\textbackslash{}Domains với value bằng 2), hạ cấp giao thức HTTP/HTTPS về Local Computer Zone (ZoneMap\textbackslash{}ProtocolDefaults với value bằng 0), chèn thẻ script độc hại, bật Proxy (ProxyEnable = 1) và thay đổi số lượng kết nối tối đa tới server (MaxConnectionsPerServer).
\item \textbf{Windows Registry Modification to Impair Security Controls or Visibility}. Phát hiện các hành vi chỉnh sửa Registry qua dòng lệnh nhằm vô hiệu hóa các cơ chế bảo mật hoặc suy giảm khả năng giám sát hệ thống; bao gồm: tắt các tính năng ghi log của PowerShell (ModuleLogging, ScriptBlockLogging, Transcription), tắt Command Prompt (DisableCMD), tắt Notification Center/Toast Notifications, tắt thông báo nâng cao của Windows Defender và vô hiệu hóa tính năng Windows Prefetch (EnablePrefetcher). Rule phát hiện được cả registry hive HKCU thay vì chỉ HKLM như PowerShell Script Block Logging Disabled
\item \textbf{Windows Explorer Functionality Restriction}. Phát hiện các hành vi chỉnh sửa khóa Registry (Software\textbackslash{}Microsoft\textbackslash{}Windows\textbackslash{}CurrentVersion\textbackslash{}Policies\textbackslash{}Explorer) để hạn chế các tính năng giao diện và quản trị của Windows Explorer; bao gồm: ẩn/vô hiệu hóa nút LogOff, ẩn Desktop (NoDesktop), ẩn/tắt tính năng Run (NoRun), Find (NoFind), Control Panel (NoControlPanel), ẩn các biểu tượng khay hệ thống (Pin, Âm thanh, Mạng, Clock) và chặn thực thi ứng dụng chỉ định (DisallowRun).
\item \textbf{Windows Policy Registry Modification}. Phát hiện các hành vi chỉnh sửa các chính sách hệ thống trong Registry (Policies\textbackslash{}System) nhằm ngăn chặn người dùng truy cập các công cụ quản trị; bao gồm: vô hiệu hóa Task Manager (DisableTaskMgr), vô hiệu hóa Registry Editor (DisableRegistryTools) và tắt tính năng thay đổi mật khẩu (DisableChangePassword).
\item \textbf{Change PowerShell Policies to an Insecure Level}. Phát hiện hành vi thay đổi chính sách powershell nhằm thực thi mã độc thông qua registry value ExecutionPolicy với data là Restricted
\item \textbf{Add SafeBoot Keys Via Reg Utility}. Phát hiện hành vi đến cấu hình khởi động an toàn SafeBoot cho phép mã độc tiếp tục chạy trong chế độ Safe Mode thông qua hành vi chỉnh sửa hoặc tạo mới các khóa Registry HKLM\textbackslash{}SYSTEM\textbackslash{}CurrentControlSet\textbackslash{}Control\textbackslash{}SafeBoot*.
\item \textbf{Modify Registry related to hermeticwiper}. Phát hiện chuỗi hành vi chỉnh sửa hai giá trị Registry ShowInfoTip và ShowCompColor về 0 tại đường dẫn Software\textbackslash{}Microsoft\textbackslash{}Windows\textbackslash{}CurrentVersion\textbackslash{}Explorer\textbackslash{}Advanced trong khoảng thời gian ngắn thông qua tiến trình reg.exe. Kỹ thuật này mô phỏng hành vi của mã độc phá hoại HermeticWiper nhằm ẩn thông tin gợi ý và màu sắc nhận diện các tệp tin bị nén/mã hóa trên Windows Explorer, giúp che giấu các thao tác can thiệp hệ thống.
\item \textbf{Microsoft Windows Defender Tampering V2}. Phát hiện hành vi can thiệp và vô hiệu hóa các tính năng bảo mật của Microsoft Defender thông qua việc chỉnh sửa các giá trị Registry. Dựa trên logic của phiên bản mặc định là Microsoft Windows Defender Tampering, rule bổ sung phát hiện chỉnh sửa registry value DisableNotifications và Notification\_Suppress, nhắm phát hiện thêm các bài Atomic Test T1112 test 38 và 51 vốn bị bỏ sót ở rule mặc định do phạm vi kiểm tra quá hẹp.
\item \textbf{Enable RDP via Remote Registry}. Phát hiện các hành vi thay đổi cấu hình Remote Desktop trong Registry nhằm cho phép kết nối trợ giúp từ xa hoặc hạ thấp cơ chế bảo mật của phiên kết nối; bao gồm: bật tính năng RDP Remote Assistance (fAllowToGetHelp = 1) và hạ mức lớp bảo mật RDP xuống mức không mã hóa (WinStations\textbackslash{}RDP-Tcp\textbackslash{}SecurityLayer = 0) nhằm duy trì quyền truy cập từ xa vào máy nạn nhân.
\item \textbf{Ursnif Malware Registry Key Creation}. Phát hiện hành vi của mã độc Ursnif download module vào lưu trong registry bằng việc phát hiện chuỗi định danh dạng GUID bằng regex nằm trong đường dẫn registry Software\textbackslash{}AppDataLow\textbackslash{}Software\textbackslash{}Microsoft\textbackslash{}".
\item \textbf{Terminal Server Client Connection History Cleared}. Phát hiện hành vi xóa dữ liệu lịch sử kết nối RDP trên hệ thống nhằm che giấu hoạt động tại khóa Registry Software\textbackslash{}Microsoft\textbackslash{}Terminal Server Client\textbackslash{}Servers hoặc Default thông qua tiến trình reg delete.
\item \textbf{RestrictedAdminMode Registry Value Tampering}. Phát hiện kẻ tấn công bật chế độ Restricted Admin Mode nhằm thực hiện kỹ thuật tấn công Pass-the-Hash qua kết nối RDP mà không cần mật khẩu bằng cách phát hiện hiện hành vi can thiệp khóa Registry (DisableRestrictedAdmin = 0) tại đường dẫn SYSTEM\textbackslash{}CurrentControlSet\textbackslash{}Control\textbackslash{}Lsa.
\item \textbf{Mimic Ransomware-Enable Multiple User Sessions}. Phát hiện hành vi bật tính năng cho phép mở nhiều phiên đăng nhập của mã độc Mimic Ransomware giúp kẻ tấn công duy trì kết nối từ xa ngầm mà không làm ngắt phiên làm việc của người dùng hợp lệ thông qua hành vi can thiệp Registry value AllowMultipleTSSessions tại đường dẫn Software\textbackslash{}Microsoft\textbackslash{}Windows\textbackslash{}CurrentVersion\textbackslash{}Winlogon
\item \textbf{RDP Registry Modification}. Phát hiện thao tác của kẻ tấn công nhằm bật/tắt dịch vụ RDP hoặc cho phép mở nhiều phiên đăng nhập đồng thời với mỗi user qua hành vi can thiệp và chỉnh sửa các khóa Registry Remote Desktop tại registry path HKLM\textbackslash{}System\textbackslash{}CurrentControlSet\textbackslash{}Control\textbackslash{}Terminal Server\textbackslash{}*.
\item \textbf{RedLine Registry Modification}. Phát hiện các thao tác can thiệp Registry mô phỏng mã độc RedLine Stealer nhằm vô hiệu hóa tính năng tự động cập nhật hệ thống Windows qua các registry value NoAutoUpdate, AUOptions, NoAutoRebootWithLoggedOnUsers, DoNotConnectToWindowsUpdateInternetLocations hoặc can thiệp giá trị TamperProtection tại đường dẫn SOFTWARE\textbackslash{}Microsoft\textbackslash{}Windows Defender\textbackslash{}Features
\item \textbf{Disable Secondary Authentication}. Phát hiện hành vi vô hiệu hóa cơ chế xác thực hai yếu tố hoặc thiết bị phụ Secondary Authentication Factor và FIDO Sign-on thông qua việc đổi các giá trị AllowSecondaryAuthenticationDevice hoặc AllowExternalDeviceSignon về 0 tại các đường dẫn Registry SOFTWARE\textbackslash{}Policies\textbackslash{}Microsoft\textbackslash{}SecondaryAuthenticationFactor và SOFTWARE\textbackslash{}Policies\textbackslash{}Microsoft\textbackslash{}FIDO nhằm suy giảm rào cản bảo mật hệ thống.
\item \textbf{Darkgate Registry Modification}. Phát hiện các thao tác can thiệp Registry đặc trưng của mã độc DarkGate bao gồm vô hiệu hóa tính năng Anti-Alias và thiết lập an toàn của Terminal Services (DisableRemoteDesktopAntiAlias = 1, DisableSecuritySettings = 1), ẩn giao diện báo lỗi hệ thống Windows Error Reporting (DontShowUI) và ghi đè cấp độ xác thực RDP Client (AuthenticationLevelOverride) nhằm suy giảm cơ chế bảo mật và hỗ trợ truy cập từ xa ngầm.
\item \textbf{Modification of IE Registry Settings}. Phát hiện hành vi chỉnh sửa hoặc thiết lập cấu hình máy chủ Proxy trong Registry tại đường dẫn Software\textbackslash{}Microsoft\textbackslash{}Windows\textbackslash{}CurrentVersion\textbackslash{}Internet Settings với giá trị ProxyServer được các dòng mã độc như DarkGate lợi dụng nhằm điều hướng lưu lượng mạng, phục vụ kết nối C2.
\item \textbf{ShimCache Flush}. Phát hiện hành vi xóa bộ nhớ đệm ShimCache bằng cách giám sát tiến trình rundll32.exe gọi thực thi các hàm xóa cache trong thư viện hệ thống apphelp.dll với ShimFlushCache/\#250 hoặc kernel32.dll với BaseFlushAppcompatCache/\#46 qua sự kiện tạo tiến trình Sysmon Event ID 1 hoặc Windows Event ID 4688. Kỹ thuật này được kẻ tấn công lợi dụng để xóa bớt dấu vết thực thi ứng dụng/tệp tin độc hại, gây cản trở cho quá trình điều tra sự cố số.
\item \textbf{Modification of BitLocker FVE Policy via Registry}. Phát hiện các hành vi can thiệp hoặc thay đổi chính sách mã hóa ổ đĩa BitLocker tại đường dẫn SOFTWARE\textbackslash{}Policies\textbackslash{}Microsoft\textbackslash{}FVE thường bị kẻ tấn công lợi dụng nhằm làm suy giảm hoặc thay đổi các cơ chế xác thực cấu hình bảo mật ổ đĩa hệ thống.
\item \textbf{Abusing Windows TelemetryController for Persistence}. Phát hiện hành vi lạm dụng khóa Registry của Windows TelemetryController tại đường dẫn SOFTWARE\textbackslash{}Microsoft\textbackslash{}Windows NT\textbackslash{}CurrentVersion\textbackslash{}AppCompatFlags\textbackslash{}TelemetryController thường bị lợi dụng để đăng ký thực thi tệp tin hoặc câu lệnh độc hại ngầm mỗi khi tác vụ đo đạc dữ liệu từ xa được khởi chạy nhằm phục vụ mục đích duy trì quyền truy cập.
\item \textbf{Abusing MyComputer Disk for Persistence}. Phát hiện hành vi thay thế đường dẫn thực thi tiện ích hệ thống cleanuppath, DefragPath, BackupPath để ghi đè bằng mã độc hại tại khóa Registry SOFTWARE\textbackslash{}Microsoft\textbackslash{}Windows\textbackslash{}CurrentVersion\textbackslash{}Explorer\textbackslash{}MyComputer được sinh ra từ tiến trình bất thường
\item \textbf{App Paths hijacking}. Phát hiện hành vi chiếm đoạt đường dẫn thực thi ứng dụng bằng cách ghi đè đường dẫn của ứng dụng hợp lệ như msedge.exe bằng tệp thực thi khác nhằm tự động khởi chạy mã độc khi người dùng mở ứng dụng. Rule sử dụng truy vấn ESQL để trích xuất và so sánh tên ứng dụng được đăng ký với tên tệp thực thi đích tại khóa Registry SOFTWARE\textbackslash{}Microsoft\textbackslash{}Windows\textbackslash{}CurrentVersion\textbackslash{}App Paths và cảnh báo khi hai giá trị này không trùng khớp.
\item \textbf{Code Executed Via Office Add-in XLL File}. Phát hiện hành vi nạp và thực thi tệp Add-in như tệp .xll độc hại trên các ứng dụng Microsoft Office như EXCEL.EXE, WINWORD.EXE... thông qua sự kiện Sysmon Event ID 7 Image/DLL Loaded. Rule kiểm tra các file DLL/Add-in được nạp từ các đường dẫn nguy cơ cao và bổ sung thêm các file không có chữ ký số tin cậy so với Prebuilt Rule Suspicious Execution via Microsoft Office Add-Ins mặc định
\item \textbf{Web Shell}. Phát hiện hành vi khởi tạo tệp tin Web Shell trên các thư mục gốc của dịch vụ máy chủ web (IIS, Apache, Nginx,...) thông qua sự kiện tạo file Sysmon Event ID 11. Rule kiểm tra các thao tác tạo file mang phần mở rộng như asp, ashx, asmx,... tại các đường dẫn lưu trữ web nhạy cảm như \textbackslash{}inetpub\textbackslash{}wwwroot\textbackslash{}, \textbackslash{}htdocs\textbackslash{}, \textbackslash{}www\textbackslash{}, \textbackslash{}html\textbackslash{}. Rule có điều kiện phát hiện rộng hơn so với Prebuilt Rule Potential IIS Web Shell File Creation, đồng thời phát hiện sớm hành vi web shell so với Web Shell Detection: Script Process Child of Common Web Processes chỉ phát hiện khi tiến trình độc hại đã được thực thi
\item \textbf{Modify Terminal Services DLL Path}. Phát hiện hành vi thay thế thư viện mặc định bằng DLL độc hại duy trì quyền truy cập mỗi khi dịch vụ RDP khởi chạy bằng cách phát hiện thay đổi đường dẫn thư viện DLL (ngoại trừ file DLL nằm trong C:\textbackslash{}Windows\textbackslash{}system32\textbackslash{}) của dịch vụ Terminal Services tại khóa Registry System\textbackslash{}CurrentControlSet\textbackslash{}services\textbackslash{}TermService\textbackslash{}Parameters với registry value ServiceDll.
\item \textbf{Add Port Monitor persistence in Registry}. Phát hiện lợi dụng cơ chế Port Monitor để tự động thực thi tệp DLL độc hại dưới tiến trình spoolsv.exe mỗi khi hệ thống khởi động bằng cách phát hiện hành vi đăng ký hoặc thêm mới file DLL vào khóa quản lý Port Monitor của dịch vụ in ấn tại đường dẫn Registry SYSTEM\textbackslash{}CurrentControlSet\textbackslash{}Control\textbackslash{}Print\textbackslash{}Monitors\textbackslash{}*\textbackslash{}Driver và không loại trừ UserId S-1-5-18 như Prebuilt Rule Potential Port Monitor or Print Processor Registration Abuse
\item \textbf{Add persistence via\_Recycle bin}. Phát hiện hành vi lợi dụng Recycle Bin để tự động thực thi mã độc mỗi khi người dùng thao tác mở thùng rác nhằm duy trì quyền truy cập, thông qua việc theo dõi các sự kiện registry ghi đè giá trị command tại khóa CLSID 645FF040-5081-101B-9F08-00AA002F954E đặc trưng của Recycle Bin
\item \textbf{Change Default File Association Via Assoc}. Phát hiện hành vi thay đổi File Association nhằm tự động thực thi mã độc khi người dùng mở các định dạng tệp tương ứng để duy trì quyền truy cập Persistence, thông qua việc theo dõi các cmd.exe có assoc trên dòng lệnh
\item \textbf{Component Object Model Hijacking v2}. Phát hiện hành vi COM Object Hijacking nhằm tự động nạp và thực thi tệp DLL độc hại khi gọi COM object tương ứng để duy trì quyền truy cập, thông qua việc giám sát thay đổi Registry path SOFTWARE\textbackslash{}Classes\textbackslash{}CLSID\textbackslash{}*\textbackslash{}InprocServer32 qua log tiến trình thay vì log registry như Prebuilt Rule Component Object Model Hijacking
\item \textbf{Custom debugger for WER}. Phát hiện hành vi lạm dụng dịch vụ báo lỗi Windows Error Reporting bằng cách đăng ký trình gỡ lỗi tùy chỉnh Custom Debugger nhằm tự động thực thi mã độc mỗi khi có ứng dụng bị treo, thông qua việc giám sát Registry path Software\textbackslash{}Microsoft\textbackslash{}Windows\textbackslash{}Windows Error Reporting\textbackslash{}Hangs và value Debugger.
\item \textbf{Event Viewer Registry Modification}. Phát hiện hành vi lạm dụng cơ chế chuyển hướng của Event Viewer nhằm tự động thực thi câu lệnh hoặc tệp tin độc hại thông qua việc giám sát tiến trình thao tác thiết lập hoặc thay đổi value MicrosoftRedirectionProgram hoặc MicrosoftRedirectionURL tại registry path Registry SOFTWARE\textbackslash{}Microsoft\textbackslash{}Windows NT\textbackslash{}CurrentVersion\textbackslash{}Event Viewer.
\item \textbf{Leverage Virtual Channels to execute custom DLL}. Phát hiện hành vi lạm dụng RDP Virtual Channels nhằm tự động nạp và thực thi tệp DLL độc hại mỗi khi kết nối RDP được khởi tạo để duy trì quyền truy cập thông qua việc giám sát thao tác đăng ký Addin tại Registry path Software\textbackslash{}Microsoft\textbackslash{}Terminal Server Client\textbackslash{}Default\textbackslash{}Addins.
\item \textbf{Modification of Boot Verification Program registry}. Phát hiện hành vi lạm dụng tính năng Boot Verification Program của hệ thống Windows nhằm tự động thực thi ứng dụng độc hại khi máy tính khởi động thành công bằng việc giám sát các sự kiện thay đổi Registry thiết lập hoặc thay đổi giá trị ImagePath tại đường dẫn System\textbackslash{}CurrentControlSet\textbackslash{}Control\textbackslash{}BootVerificationProgram.
\item \textbf{New Custom Shim Database Created}. Phát hiện hành vi lạm dụng cơ chế Application Shimming nhằm duy trì quyền truy cập thông qua việc giám sát sự kiện khởi tạo file Sysmon Event ID 11 thư mục hệ thống C:\textbackslash{}Windows\textbackslash{}apppatch\textbackslash{}Custom*.
\item \textbf{Persistence using STARTUP-PATH in MS-WORD}. Phát hiện hành vi lạm dụng cấu hình thư mục khởi động của Microsoft Word nhằm tự động nạp và thực thi các file hại mỗi khi ứng dụng MS Word khởi chạy để duy trì quyền truy cập, thông qua việc giám sát thao tác thiết lập registry value STARTUP-PATH tại Registry path Software\textbackslash{}Microsoft\textbackslash{}Office\textbackslash{}*Word\textbackslash{}Options*.
\item \textbf{Persistence via Active Setup}. Phát hiện hành vi lạm dụng cơ chế Active Setup của Windows nhằm tự động thực thi tệp tin hoặc câu lệnh độc hại mỗi khi người dùng đăng nhập hệ thống bằng việc giám sát tiến trình thao tác tạo hoặc chỉnh sửa thuộc tính StubPath tại đường dẫn Registry SOFTWARE\textbackslash{}Microsoft\textbackslash{}Active Setup\textbackslash{}Installed Components.
\item \textbf{Persistence via Windows Context Menu}. Phát hiện hành vi lạm dụng menu chuột phải của Windows nhằm tự động thực thi tệp tin hoặc kịch bản độc hại khi người dùng nhấp chuột phải thông qua việc giám sát thao tác đăng ký có phần mở rộng là exe, dll, bat, cmd, ps1, vbs, js tại đường dẫn Registry Directory\textbackslash{}Background\textbackslash{}shell*\textbackslash{}command.
\item \textbf{Persistence with Custom DLL}. Phát hiện hành vi đăng ký các tệp DLL tùy chỉnh độc hại nhằm tự động nạp DLL khi các tiến trình kết nối Internet, trong khi thực hiện phiên RDP thông qua việc giám sát thao tác thiết lập đường dẫn tệp tin .dll tại các khóa Registry tương ứng.
\item \textbf{Potential Suspicious Activity Using SeCEdit}. Phát hiện hành vi lạm dụng công cụ secedit.exe nhằm tạo Run Key với registry hive HKLM bằng việc phát hiện các tiến trình thực thi secedit.exe đi kèm các tham số dòng lệnh cấu hình hoặc nhập cơ sở dữ liệu chính sách /configure, /db, /import, /cfg. Run Key do secedit.exe tạo ra không bị phát hiện bởi Prebuilt Rule Startup or Run Key Registry Modification do có user SID là S-1-5-18 (SYSTEM)
\item \textbf{Powershell Execute COM Object}. Phát hiện sử dụng PowerShell để khởi tạo và thực thi COM Object thông qua mã định danh CLSID nhằm thực hiện Component Object Model Hijacking. Rule phát hiện PowerShell Script Block Logging (Event ID 4104) khi ghi nhận mã lệnh chứa đoạn truy vấn CreateInstance([type]::GetTypeFromCLSID) tương ứng với khởi tạo một đối tượng COM.
\item \textbf{Potential Remote Desktop Shadowing Activity v2}. Dựa trên logic của rule mặc định Potential Remote Desktop Shadowing Activity và thực hiện ánh xạ thêm technique T1112
\item \textbf{Registry key creation and/or modification events for SDB}. Phát hiện hành vi lạm dụng cơ chế Application Shimming qua việc giám sát tạo mới hoặc thay đổi khóa đăng ký Shim Database tùy chỉnh tại các đường dẫn Registry SOFTWARE\textbackslash{}Microsoft\textbackslash{}Windows NT\textbackslash{}CurrentVersion\textbackslash{}AppCompatFlags\textbackslash{}InstalledSDB* và SOFTWARE\textbackslash{}Microsoft\textbackslash{}Windows NT\textbackslash{}CurrentVersion\textbackslash{}AppCompatFlags\textbackslash{}Custom\textbackslash{}*.
\item \textbf{Registry to load Arbitrary DLL into LSASS}. Phát hiện hành vi nạp thư viện DLL tùy chỉnh độc hại vào tiến trình lsass.exe bằng việc giám sát các sự kiện thay đổi Registry tại các registry path CurrentControlSet\textbackslash{}Services\textbackslash{}NTDS\textbackslash{}DirectoryServiceExtPt và CurrentControlSet\textbackslash{}Services\textbackslash{}NTDS\textbackslash{}LsaDbExtPt
\item \textbf{Shortcut Modification}. Phát hiện hành vi tạo hoặc sửa đổi các phím tắt có thể thực thi chương trình trong quá trình khởi động hệ thống hoặc đăng nhập của người dùng bằng cách giám sát sự kiện tạo file có đuôi là .lnk, .url hoặc nằm trong các thư mục Startup và được thực hiện bởi các tiến trình dòng lệnh bất thường như cmd, PowerShell, mshta,... hoặc ứng dụng Microsoft Office như WINWORD, EXCEL...
\item \textbf{Suspicious Copy From or To System Directory}. Phát hiện hành vi sao chép đáng ngờ các tệp tin hoặc chương trình từ/đến các thư mục hệ thống nhạy cảm như System32, SysWOW64, WinSxS bằng các tiến trình thực thi lệnh sao chép thông qua các công cụ dòng lệnh như cmd, PowerShell, robocopy, xcopy...,
\item \textbf{Suspicious Driver Install by pnputil.exe}. Phát hiện hành vi cài đặt driver độc hại hoặc không rõ nguồn gốc vào hệ thống bằng các sự kiện khởi tạo tiến trình pnputil.exe đi kèm các tham số dòng lệnh thêm hoặc cài đặt driver (-i, install, -a, add-driver, .inf).
\item \textbf{Suspicious Modify Registry of Current User Profile}. Phát hiện hành vi ẩn phần mở rộng của tệp tin che giấu các tệp tin độc hại để duy trì quyền truy cập qua sự kiện sửa đổi giá trị HideFileExt tại đường dẫn Registry Software\textbackslash{}Microsoft\textbackslash{}Windows\textbackslash{}CurrentVersion\textbackslash{}Explorer\textbackslash{}Advanced.
\item \textbf{UEFI Persistence Via Wpbbin}. Trong quá trình khởi động, tiến trình smss.exe tìm cấu trúc chứa nhãn Windows Platform Binary Table trong bộ nhớ UEFI và tự động trích xuất mã nhị phân đó vào C:\textbackslash{}Windows\textbackslash{}System32\textbackslash{}wpbbin.exe và thực thi với quyền cao nhất. Rule phát hiện hành vi nạp mã độc vào firmware UEFI thông qua hành vi tạo file C:\textbackslash{}Windows\textbackslash{}System32\textbackslash{}wpbbin.exe
\item \textbf{Windows Accessibility Features Registry Modification}. Phát hiện hành vi lạm dụng Assistive Technologies như công cụ hỗ trợ người khuyết tật nhằm thực thi mã độc qua việc chỉnh Registry value StartEXE tại Registry path Software\textbackslash{}Microsoft\textbackslash{}Windows NT\textbackslash{}CurrentVersion\textbackslash{}Accessibility\textbackslash{}ATs*.
\item \textbf{Winlogon Userinit Key Persistence}. Phát hiện hành vi lạm dụng cơ chế Winlogon Helper DLL nhằm tự động thực thi tệp tin độc hại mỗi khi người dùng đăng nhập hệ thống qua việc thao tác chỉnh sửa hoặc thiết lập các registry value Userinit hoặc Notify tại registry path Registry SOFTWARE\textbackslash{}Microsoft\textbackslash{}Windows NT\textbackslash{}CurrentVersion\textbackslash{}Winlogon.
\item \textbf{Password Filter DLL}. Phát hiện hành vi can thiệp vào quy trình xác thực của hệ thống bằng cách đăng ký thư viện DLL tùy chỉnh nhằm đánh cắp mật khẩu hoặc duy trì quyền truy cập bằng các sự kiện thay đổi Registry tại registry value Notification Packages thuộc registry path Registry SYSTEM\textbackslash{}*ControlSet*\textbackslash{}Control\textbackslash{}Lsa.
\item \textbf{Python Site or User Customize File Creation Win}. Phát hiện hành vi Python Startup Hooks nhằm tự động thực thi mã độc mỗi khi môi trường Python được kích hoạt bằng các sự kiện tạo file Sysmon Event ID 11 sitecustomize.py hoặc usercustomize.py tại các đường dẫn thư mục site-packages của Python.
\item \textbf{Image File Execution Options Injection V2}. Phát hiện hành vi lạm dụng tính năng Image File Execution Options vốn dùng để gỡ lỗi và theo dõi tiến trình nhằm thực thi mã độc hại. Logic phát hiện tương tự rule mặc định Image File Execution Options Injection và bổ sung thêm ánh xạ đến technique T1546.008
\item \textbf{Uncommon Registry Persistence Change v2}. Phát hiện hành vi bất thường thay đổi Registry ít phổ biến. Logic phát hiện tương tự rule mặc định Uncommon Registry Persistence Change và bổ sung thêm ánh xạ đến technique T1547.004, T1037.001
\item \textbf{Suspicious Service was Installed in the System v2}. Phát hiện hành vi khởi tạo hoặc đăng ký một dịch vụ bất thường hoặc không thuộc danh sách tin cậy thông qua việc Event ID 4697 trong Security Log hoặc Event ID 7045 trong System Log và loại trừ một số dịch vụ hợp lệ chuẩn của hệ thống và phần mềm quản trị.
\item \textbf{Suspicious Windows Powershell Arguments v2}. Phát hiện PowerShell khởi chạy với tham số hoặc chuỗi lệnh đáng ngờ như mã hóa, tải nội dung từ mạng, giải mã, phản chiếu .NET hoặc kỹ thuật làm rối lệnh.
\end{enumerate}
